\begin{center}
Problema C

Halloween no Bairro

{\small Nome do arquivo: C.c, C.cpp, C.java, ou C.py}

{\large Timelimit: $1$}

{\tiny Por André da Cunha Ribeiro}

\end{center}			
			
Neste Halloween do Abóbora Contest $5$, várias crianças decidiram sair para pedir doces em seu bairro. Cada criança tem um número de identificação de $1$ a $N$, e todas as casas no bairro distribuem doces em uma ordem aleatória. Ao final da noite, cada criança retorna para contar quantos doces coletou e, com base nesses números, você deve determinar a posição de cada criança no ranking de doces recebidos.

Por exemplo, se $N = 5$ e as quantidades de doces coletados pelas crianças foram: $8$, $15$, $7$, $20$, $10$, então a criança com número $1$ ficou na $4$ª posição, a criança com número $2$ ficou na $2$ª posição, a criança com número $3$ ficou na $5$ª posição, a criança com número $4$ ficou na $1$ª posição e a criança com número $5$ ficou na $3$ª posição.

Sua tarefa é ajudar a classificar as crianças de acordo com a quantidade de doces recebidos, da seguinte forma: a criança que coletou mais doces deve ficar na $1$ª posição, a que coletou o segundo maior número de doces fica na $2$ª posição, e assim por diante.

\vspace{0.3cm}
\textbf{Entrada}
A primeira linha da entrada contém um único inteiro $N$, representando o número de crianças que participaram. Nas próxima linha contêm $N$ inteiro, representando a quantidade de doces coletados pela criança de número correspondente.
\begin{itemize}
    \item Cada criança coleta uma quantidade diferente de doces.
\end{itemize}
\begin{itemize}
\item $1 \leqslant |N| \leqslant 10.000$;
\end{itemize}

\vspace{0.3cm}
\textbf{Saída}
Seu programa deve imprimir um linha contendo os $N$ inteiro separados por um espaço. A $i-$ésima posição deve conter a posição no ranking da criança com número $i$. O último valor não tem espaço em branco, somente o terminador de linha.

\begin{table}[h]
    \centering
    \begin{tabular}{|p{7cm}|p{7cm}|}
        \hline
        \begin{center}
            \vspace{-0.3cm}
            \textbf{Exemplos de Entrada}
            \vspace{-0.3cm}
        \end{center} 
         &  
        \begin{center}
            \vspace{-0.3cm}
            \textbf{Exemplos de Saída}
            \vspace{-0.3cm}
        \end{center} \\ \hline
        \texttt{5} & \texttt{4 2 5 1 3} \\ 
        \texttt{8 15 7 20 10} & \texttt{} \\ \hline
        \texttt{3} & \texttt{1 3 2} \\ 
        \texttt{12 5 8} & \texttt{} \\ \hline
    \end{tabular}
    \caption{Exemplos de entradas e saídas}
    \label{tab:inputc3}
\end{table}

\newpage

\begin{center}
\noindent
\lstinputlisting{abobora_24/C/C4.cpp}
\end{center}